%! Bias in Samples and Surveys — Unit 1, Lesson 1.4 — Guided Notes & Practice (Answer Key)
% THE in-class document, in the MAIN IDEAS / NOTES shape (LESSON_SHAPE.md §1, ported from AP
% Statistics 2026-09-12): the vocabulary box, the hook, then ONE two-column guidednotes table.
% Each row is one idea — a label on the left; on the right one or two COMPLETE printed
% sentences, ONE large pre-drawn display the student reads or names, then a bold prompt and a
% two-across grid of numbered problems with 1.0–1.4cm of writing room. Problems run 1..19.
% DENSITY: a \blank{} only where a word IS the answer (the five cells of the diagnostic table
% in row 3 — the one \wordbank strip sits above it); no mid-sentence blanks; every display is
% pre-drawn; the page belongs to the student's pen.
%   I DO   : the teacher reads the definition, marks up the display, works the first problem
%            of each grid (1, 5, 10, 12).
%   WE DO  : the class works the rest of each grid, students holding the pen; the last row,
%            Guided Practice (16–19), is worked entirely together on Cedar Ridge Apartments
%            instead of the pool.
%   YOU DO : nothing on this page — ../homework/ IS the individual practice, started in the
%            last ten minutes of the period. No independent set, no reflection box, no
%            objectives box (the cover carries the targets).
% THE TRAP is problem 8 (row 2: all 1,500 got a card, 200 mailed it back — a perfect frame and
% a broken study: nonresponse, not undercoverage). Problem 4 is the smaller trap (the tighter
% Tuesday row is the wrong one — consistent is not correct).
% THE CRUX is problem 13 (row 4, "Bigger Is Not Better"): ten Tuesdays, 270/750 = 36%, the
% same 16-point miss — a larger sample is more PRECISE, not more CORRECT; size shrinks
% variability and does nothing to direction (PS.DC.2c, AFDA.DA.2e). The hook vote is re-taken
% there. Problem 18 is the crux transferred to Cedar Ridge; problem 19 is AFDA.DA.2g (name the
% bias in the wording, rewrite it balanced).
%
% THE DATA
%   Harbor Point Community Pool (warm-up + notes): 1,500 members; census truth 300/1500 = 20%
%     mainly swim laps. Tuesday-evening sample 27/75 = 36% -> 540 predicted, off by 240.
%     Three honest samples 19/22/19 average 20%; three Tuesday samples 35/36/37 average 36%.
%     Ten Tuesdays: 270/750 = 36% — the same 16-point miss. Question order: 62% vs 41% = 21.
%   Guided Practice, Cedar Ridge Apartments: 600 households ALL mailed a card; 150 return
%     (25%); 96/150 = 64% -> 384 predicted; census 240/600 = 40%, so 24 points and 144
%     households too many; re-run 192/300 = 64% again.
%
% PAGE LOCKSTEP with notes_key/: the key is GENERATED from this file by templates/lesson/mkkey.py
% (spec: templates/lesson/examples/lesson04_notes_key.json) — every \blank{} becomes an \ans{},
% every \vterm{} a \vtermans{}{}, every \par\writelines{n} n \ansline{} calls, and the answer
% argument of every \pcell, \writespace and \labelbox is filled. Answer spaces are fixed
% heights, so the two files cannot drift. KEEP EVERY \pcell ON ONE SOURCE LINE — mkkey matches
% line by line.
% TABLE MECHANICS: inside a cell, `\\` ENDS THE ROW — break lines with \par. A row cannot break
% across pages: the figure gets its own sub-row (`\\` then `& ...`) and EACH GRID ROW is its own
% sub-row — close the probgrid, `\\ &`, reopen as probgrid* (no top rule).
\documentclass[12pt]{article}
\usepackage{saar-article}
\usepackage{saar-key}   % pulls in -boxes; do NOT also load -boxes
\newcommand{\TallMath}[1]{$\displaystyle #1\rule[-1.4em]{0pt}{3.2em}$}
% Word bank strip — ONLY above a table the student fills (row 3). Defined per document;
% the key inherits it and prints the same strip, so heights match.
\newcommand{\wordbank}[1]{\par\vspace{2pt}\noindent\fcolorbox{goldacc}{goldbg}{\parbox{\dimexpr\linewidth-2\fboxsep-2\fboxrule\relax}{\small\textbf{Word bank:} #1}}\par\vspace{4pt}}
% Vocabulary rows: a FIXED-HEIGHT row carrying the term label and OPEN writing space —
% no underline beside the term and no rule line (user correction 2026-09-06: the black
% answer line crowded the row; students write beside and below the term). The key fills
% exactly the same height, so the box cannot drift blank vs. keyed. saar-article's
% \termblank adds an inline blank and a rule line, so the pair is defined here (the key is
% generated from this file and inherits both). 2.3cm per row gives students three
% handwritten lines of room; keep key definitions to two lines.
\newcommand{\termrowheight}{2.3cm}
\newlength{\termrowinset}\setlength{\termrowinset}{7pt}
\newcommand{\vterm}[1]{%
  \par\noindent\begin{minipage}[t][\termrowheight][t]{\linewidth}%
    \vspace*{\termrowinset}%
    \textbf{\textcolor{cerulean}{#1:}}%
  \end{minipage}\par}
\newcommand{\vtermans}[2]{%
  \par\noindent\begin{minipage}[t][\termrowheight][t]{\linewidth}%
    \vspace*{\termrowinset}%
    \textbf{\textcolor{cerulean}{#1:}}\quad\ans{#2}%
  \end{minipage}\par}

\begin{document}

\pageheader{Unit 1, Lesson 1.4}{Guided Notes \& Practice --- Answer Key}
% NAMESTRIP: no \namedateperiod here — mirrors the blank. NO teachernote here —
% teacher prose lives in the lesson plan.

\vspace{0.15in}

\begin{vocabbox}
\small
Fill in each term \textbf{as we name it} in the notes below.
\vtermans{Bias}{A study is biased when its procedure systematically misses the population value the same way every time --- a direction, not bad luck.}
\vtermans{Sampling bias}{Bias from how the sample was chosen: some members of the population are more likely to be picked than others.}
\vtermans{Undercoverage}{Part of the population is left off the sampling frame, so those members have no chance at all of being chosen.}
\vtermans{Nonresponse}{Members are chosen properly but never answer, so the data is missing exactly the people who stayed silent.}
\vtermans{Response bias}{Anything about the survey itself --- who asks, the wording, the order --- that pushes a respondent off their true answer.}
\end{vocabbox}

\boxguard[12]
\begin{hookbox}
\small
Last class the board threw out the Tuesday-evening study: $27$ of the $75$ members asked said
they mainly swim laps --- $36\%$ --- while the office's own records say $20\%$. The staff
member says the problem was \emph{size}. She will run the same Tuesday-evening survey for
ten weeks and ask $750$ members instead of $75$.

\vspace{3pt}
\textbf{Will her new number land closer to the truth, $20\%$?} Circle one: \quad yes \qquad
no \qquad it depends \quad --- and say why you think so. We vote again at problem~13.
\par\ansline{Answers vary --- most of the room votes yes: more people must mean a better number.}
\ansline{(Settled at problem 13: no --- $270/750 = 36\%$, the same 16-point miss.)}
\end{hookbox}

\small
\begin{guidednotes}
% ── ROW 1: bias is a direction, not bad luck ─────────────────────────────────
\mainidea[A direction, not bad luck]{Bias} &
A study is \textbf{biased} when its procedure \emph{systematically} overestimates or
underestimates the population value --- every sample leans the \textbf{same way}. Honest
random samples (Lesson 1.2) only \emph{scatter around} the truth. The records say $20\%$
mainly swim laps; three samples of $75$, two ways:
\\
& {\centering
\begin{tikzpicture}[scale=0.70]
  % percent axis: x = (p - 10) * 0.32, so 10% -> 0 and 45% -> 11.2
  \draw[->, thick] (-0.2,0) -- (11.6,0);
  \foreach \p in {10,15,...,45} {
    \draw ({(\p-10)*0.32},-0.1) -- ({(\p-10)*0.32},0.1);
    \node[below, font=\tiny] at ({(\p-10)*0.32},-0.1) {\p\%}; }
  % the truth: 20%
  \draw[cerulean, very thick, dashed] (3.2,-0.05) -- (3.2,2.45);
  \node[font=\scriptsize\bfseries\color{cerulean}, above] at (3.2,2.45) {truth: $20\%$ (the records)};
  % row A: random from the full list — 19, 22, 19
  \node[anchor=east, font=\scriptsize\bfseries] at (-0.3,1.8) {random from the list};
  \draw[linegray] (0,1.8) -- (4.5,1.8);
  \fill[goldacc] (2.88,1.67) circle (0.13);
  \fill[goldacc] (2.88,1.95) circle (0.13);
  \fill[goldacc] (3.84,1.8) circle (0.13);
  \node[font=\tiny, anchor=west] at (4.75,1.8) {$19\%$, $22\%$, $19\%$};
  % row B: whoever arrived Tuesday evening — 35, 36, 37
  \node[anchor=east, font=\scriptsize\bfseries] at (-0.3,0.95) {Tuesday evening};
  \draw[linegray] (0,0.95) -- (8.85,0.95);
  \fill[redacc] (8.0,0.95) circle (0.13);
  \fill[redacc] (8.32,0.95) circle (0.13);
  \fill[redacc] (8.64,0.95) circle (0.13);
  \node[font=\tiny, anchor=west] at (9.1,0.95) {$35\%$, $36\%$, $37\%$};
  % the gap
  \draw[<->, redacc, thick] (3.2,0.4) -- (8.32,0.4);
  \node[font=\tiny\bfseries\color{redacc}, above] at (5.76,0.42) {$16$ points, every time};
\end{tikzpicture}\par\vspace{3pt}
Scatter around it: \labelbox{2.4cm}{variability}\qquad Same miss every time: \labelbox{1.3cm}{bias}\par}
\\
& \notesprompt{Read the display. \emph{Consistent is not the same as correct.}}
\begin{probgrid}{|Y|Y|}
\pcell{1}{Average each row: random ($19$, $22$, $19$) and Tuesday ($35$, $36$, $37$).}{1.0cm}{Random $20\%$; Tuesday $36\%$.} & \pcell{2}{Both rows wobble. Which row misses the truth \emph{every} time --- too high or too low?}{1.0cm}{The Tuesday row --- always too high.} \\ \hline
\end{probgrid}
\\
& \begin{probgrid}*{|Y|Y|}
\pcell{3}{Tuesday evening is lap-swim time. Who is over-counted, and why does that push the number up?}{1.2cm}{Lap swimmers: Tuesday evening is lap-swim time.} & \pcell{4}{The Tuesday row agrees with itself better. Is it the row to trust? Explain.}{1.2cm}{No --- consistent is not correct.} \\ \hline
\end{probgrid}
\\ \hline
% ── ROW 2: sampling bias — undercoverage and nonresponse ─────────────────────
\mainidea[Who never had a chance]{Sampling Bias} &
\textbf{Sampling bias} gets in when some members are more likely to be chosen than others ---
it arrives \emph{before anyone speaks}. \textbf{Undercoverage:} part of the population is left
off the \textbf{sampling frame} and can never be chosen. \textbf{Nonresponse:} members are
chosen properly but never answer.
\\
& {\centering
\begin{tikzpicture}[scale=0.70]
  % the population
  \fill[frost, rounded corners=6pt] (0,0) rectangle (13.0,3.0);
  \draw[cerulean, thick, rounded corners=6pt] (0,0) rectangle (13.0,3.0);
  \node[anchor=west, font=\scriptsize\bfseries\color{cerulean}] at (0.2,2.7) {Harbor Point --- all $1{,}500$ members};
  % the frame: who could be chosen
  \fill[white, rounded corners=4pt] (0.3,0.3) rectangle (8.6,2.3);
  \draw[cerulean, thick, rounded corners=4pt] (0.3,0.3) rectangle (8.6,2.3);
  \node[anchor=west, font=\tiny\bfseries\color{cerulean}] at (0.45,2.05) {the frame: members who \emph{could} be chosen};
  % the chosen: an oval inside the frame; filled dots answered, hollow dots silent
  \fill[goldbg] (4.45,1.15) ellipse (3.0 and 0.72);
  \draw[goldacc, very thick] (4.45,1.15) ellipse (3.0 and 0.72);
  \node[font=\tiny\bfseries\color{goldacc}, fill=goldbg, inner sep=1pt] at (4.45,1.63) {chosen};
  \foreach \x in {2.2,2.7,3.2,3.7,4.2} { \fill[slate] (\x,1.2) circle (0.1); }
  \foreach \x in {4.9,5.4,5.9,6.4,6.9} { \draw[slate, thick] (\x,1.2) circle (0.1); }
  \node[font=\tiny, anchor=north] at (3.2,1.0) {\emph{answered}};
  \node[font=\tiny, anchor=north] at (5.9,1.0) {\emph{silent}};
  % off the frame
  \foreach \x in {9.3,9.9,10.5,11.1,11.7,12.3} { \foreach \y in {0.55,1.05} { \fill[slate!50] (\x,\y) circle (0.09); }}
  \node[font=\tiny\bfseries\color{redacc}, align=center] at (10.8,1.85) {off the frame:\\could \emph{never} be chosen};
\end{tikzpicture}\par\vspace{3pt}
Off the frame: \labelbox{2.9cm}{undercoverage}\quad Chosen but silent: \labelbox{2.4cm}{nonresponse}\par}
\textbf{The test: \emph{could that person have been chosen?}} No $\Rightarrow$ undercoverage.
Yes, but silent $\Rightarrow$ nonresponse.
\\
& \notesprompt{Who is missing, and which bias? Ask the test question.}
\begin{probgrid}{|Y|Y|}
\pcell{5}{Ask the members who arrive on a Tuesday evening.}{1.0cm}{Undercoverage --- daytime members are off the frame.} & \pcell{6}{Email the survey only to members with an email address on file.}{1.0cm}{Undercoverage --- no email, no chance.} \\ \hline
\end{probgrid}
\\
& \begin{probgrid}*{|Y|Y|}
\pcell{7}{Post the survey on the pool's social media page.}{1.0cm}{Undercoverage --- non-users cannot be chosen.} & \pcell{8}{Mail a card to \emph{all} $1{,}500$; $200$ mail it back. \textbf{Careful: the frame is perfect.}}{1.2cm}{Nonresponse --- the frame was perfect.} \\ \hline
\end{probgrid}
\\ \hline
% ── ROW 3: response bias — after the right person is chosen ──────────────────
\mainidea[After the right person is chosen]{Response Bias} &
\textbf{Response bias} is anything about the survey itself --- who asks, the wording, the
order --- that pushes a respondent from their true answer. It comes \emph{after} the right
person is chosen, so no sampling method prevents it. Seven kinds:
\\
& \footnotesize\renewcommand{\arraystretch}{1.0}
\begin{tabularx}{\linewidth}{@{} >{\bfseries}p{3.5cm} X @{}}
\toprule
\textbf{\normalfont Kind} & \textbf{What the respondent does} \\
\midrule
Demand & Guesses what the study wants to hear and says it. \\
Social desirability & Gives the answer that looks good to others. \\
Dissent & Says no to everything, often without reading. \\
Acquiescence & Agrees with everything --- every answer positive. \\
Extreme responses & Always ``strongly agree'' or ``strongly disagree.'' \\
Neutral responding & Always the middle, no-opinion choice. \\
Question order & An earlier question changes a later answer. \\
\bottomrule
\end{tabularx}\small
\\
& \textbf{9.} Name the kind of response bias in each situation at the pool.
\wordbank{demand \;\textbullet\; social desirability \;\textbullet\; dissent \;\textbullet\; acquiescence \;\textbullet\; extreme responses \;\textbullet\; neutral responding \;\textbullet\; question order}
\footnotesize\renewcommand{\arraystretch}{1.05}
\begin{tabularx}{\linewidth}{@{} X p{3.4cm} @{}}
\toprule
\textbf{At the pool} & \textbf{Kind} \\
\midrule
The lifeguard hands it out and waits; the staff is rated ``excellent.'' & \ans{social desirability} \\
``No opinion'' on all fifteen questions. & \ans{neutral responding} \\
Agrees with every statement, even two that conflict. & \ans{acquiescence} \\
Titled \emph{``Help us show the pool needs longer hours.''} & \ans{demand} \\
Crowding asked first: $62\%$ back a lane. Money first: $41\%$. & \ans{question order} \\
\bottomrule
\end{tabularx}\small
\\
& \notesprompt{Say which way it pushes.}
\begin{probgrid}{|Y|Y|}
\pcell{10}{Same words, two orders: $62\%$ versus $41\%$. How many points apart --- and which order would a member who \emph{wants} the lane ask?}{1.0cm}{$21$ points; crowding first.} & \pcell{11}{Choose the $75$ at random from the full list, but keep the title \emph{``Help us show the pool needs longer hours.''} Fixed? Explain.}{1.0cm}{No --- the title is response bias.} \\ \hline
\end{probgrid}
\\ \hline
% ── ROW 4: THE CRUX — bigger is not better ───────────────────────────────────
\mainidea[Ten Tuesdays]{Bigger Is Not Better} &
Told her sample was too small, the staff member runs the \emph{same} Tuesday survey for ten
weeks: $750$ asked, and $270$ say they mainly swim laps.
\\
& {\centering
\begin{tikzpicture}[scale=0.65]
  % percent axis: x = (p - 10) * 0.32
  \draw[->, thick] (-0.2,0) -- (11.6,0);
  \foreach \p in {10,15,...,45} {
    \draw ({(\p-10)*0.32},-0.1) -- ({(\p-10)*0.32},0.1);
    \node[below, font=\tiny] at ({(\p-10)*0.32},-0.1) {\p\%}; }
  \draw[cerulean, very thick, dashed] (3.2,-0.05) -- (3.2,2.35);
  \node[font=\scriptsize\bfseries\color{cerulean}, above] at (3.2,2.35) {truth: $20\%$};
  % one Tuesday: 75 asked — a wide bracket of where it could land, centred at 36.
  % Captions sit to the RIGHT of each bracket, never above it: TikZ scales coordinates
  % but NOT font size, so a caption stacked above a line collides at scale < 0.75.
  \node[anchor=east, font=\scriptsize\bfseries] at (-0.3,2.0) {one Tuesday, $75$};
  \draw[redacc, thick, |-|] (6.4,2.0) -- (10.24,2.0);
  \fill[redacc] (8.32,2.0) circle (0.14);
  \node[font=\tiny, anchor=west] at (10.4,2.0) {wide};
  % ten Tuesdays: 750 asked — a narrow bracket, still centred at 36
  \node[anchor=east, font=\scriptsize\bfseries] at (-0.3,1.1) {ten Tuesdays, $750$};
  \draw[redacc, thick, |-|] (7.84,1.1) -- (8.8,1.1);
  \fill[redacc] (8.32,1.1) circle (0.14);
  \node[font=\tiny, anchor=west] at (8.95,1.1) {tighter --- still $36\%$};
  % the gap that does not move
  \draw[<->, redacc, thick] (3.2,0.40) -- (8.32,0.40);
  \node[font=\tiny\bfseries\color{redacc}, above] at (5.60,0.42) {$16$ points --- unchanged};
\end{tikzpicture}\par\vspace{3pt}
The bracket shrinks: \labelbox{2.4cm}{variability}\qquad The gap stays: \labelbox{1.3cm}{bias}\par}
\\
& \fcolorbox{cerulean}{frost}{\parbox{\dimexpr\linewidth-2\fboxsep-2\fboxrule\relax}{\centering
\textbf{A larger sample makes an estimate more \emph{precise}, not more \emph{correct}.}\par
Size shrinks \textbf{variability}; nothing about size touches \textbf{direction}. Only a better \emph{procedure} removes bias.}}
\\
& \notesprompt{Now settle the hook vote.}
\begin{probgrid}{|Y|Y|}
\pcell{12}{Compute $270$ of $750$ as a percent. Off by how many points from $20\%$? One Tuesday was off by how many?}{1.2cm}{$270/750 = 36\%$; off by $16$ --- so was one Tuesday.} & \pcell{13}{\textbf{Vote again:} yes \; no \; it depends. Did $750$ land closer than $75$? What did ten times the work change --- and not change?}{1.4cm}{No. Ten Tuesdays shrank the spread, not the lean.} \\ \hline
\end{probgrid}
\\
& \begin{probgrid}*{|Y|Y|}
\pcell{14}{Name the bias each fix removes: (a) a complete list as the frame; (b) calling back the silent; (c) neutral, anonymous wording.}{1.2cm}{(a) undercoverage (b) nonresponse (c) response bias} & \pcell{15}{The board offers a random sample of just $75$ from the full list. Smaller than $750$ --- is it better? Why?}{1.2cm}{Yes --- a random $75$ has no lean.} \\ \hline
\end{probgrid}
\\ \hline
% ── ROW 5: GUIDED PRACTICE — we do, together, a fresh context ────────────────
\mainidea[Guided practice]{Cedar Ridge} &
Cedar Ridge Apartments has $600$ households. Management mails a parking survey to
\emph{every} household; $150$ cards come back, $96$ of those say parking is a problem --- and
management is about to ask the owners for a bigger lot.
\\
& {\centering
\begin{tikzpicture}[scale=0.65]
  \fill[frost, rounded corners=6pt] (0,0) rectangle (13.0,2.4);
  \draw[cerulean, thick, rounded corners=6pt] (0,0) rectangle (13.0,2.4);
  \node[font=\scriptsize\bfseries\color{cerulean}] at (6.5,2.1) {all $600$ households --- \emph{every} one mailed a card};
  % the 150 who answered
  \fill[goldbg] (3.3,0.95) ellipse (3.1 and 0.72);
  \draw[goldacc, very thick] (3.3,0.95) ellipse (3.1 and 0.72);
  \node[font=\scriptsize\bfseries\color{goldacc}] at (3.3,1.22) {$150$ mailed it back};
  \node[font=\tiny] at (3.3,0.78) {$96$ say parking is a problem};
  % the silent
  \node[font=\scriptsize\bfseries\color{slate}] at (9.9,1.3) {$450$ threw the card away};
  \node[font=\tiny\itshape] at (9.9,0.75) {a census later: $240$ of $600$ complain};
\end{tikzpicture}\par}
\\
& \notesprompt{We work these together.}
\begin{probgrid}{|Y|Y|}
\pcell{16}{What share of the $600$ answered at all? What share of \emph{those} said parking is a problem? Scale that percent to all $600$.}{1.3cm}{$25\%$ answered; $64\%$ of those; $384$ of $600$.} & \pcell{17}{Every household got a card. Could a silent one have been chosen --- so which bias? The census says $40\%$: off by how many points, how many households, which way?}{1.4cm}{Yes --- nonresponse. $24$ points, $144$ households high.} \\ \hline
\end{probgrid}
\\
& \begin{probgrid}*{|Y|Y|}
\pcell{18}{Management mails again: $300$ back, $192$ complain --- $64\%$ again. Two-bedroom households own the most cars \emph{and} answer at twice the rate. Why did twice the cards not help?}{1.2cm}{Same procedure, more of it --- the same households.} & \pcell{19}{The card asked: \emph{``Don't you agree that Cedar Ridge's parking lot is dangerously overcrowded?''} Name the response bias it invites, then rewrite it to push neither way.}{1.2cm}{Acquiescence. ``Is parking a problem for you?''} \\ \hline
\end{probgrid}
\\ \hline
\end{guidednotes}

% ── THE NOTES END AT GUIDED PRACTICE ─────────────────────────────────────────
% There is deliberately no "you do" here: ../homework/ is the individual practice,
% started in class in the last ten minutes while the teacher circulates.

\end{document}
